\section{Validasi Kode: Aturan dan Multi-Pass}
\label{sec:ch16-validasi}

\subsection{Aturan Validasi}

Validasi memeriksa aturan yang tidak tercakup oleh tipe:
\begin{enumerate}[label=\arabic*.,leftmargin=*]
  \item \code{break} dan \code{continue} hanya di dalam loop;
  \item fungsi non-\code{void} harus mengembalikan nilai;
  \item \code{main} harus \code{int main()} tanpa parameter;
  \item \code{return} hanya di dalam fungsi;
  \item pernyataan setelah \code{return} tidak terjangkau;
  \item pemanggilan fungsi sesuai signature.
\end{enumerate}

Aturan dapat diperluas sesuai kebutuhan bahasa. Yang penting adalah
setiap aturan memiliki pesan kesalahan yang jelas dan lokasi yang
tepat.

\subsection{Multi-Pass}

Validasi dilakukan setelah analisis semantik:
\begin{enumerate}[label=\arabic*.,leftmargin=*]
  \item parsing menghasilkan AST;
  \item analisis semantik membangun tabel simbol dan memeriksa tipe;
  \item validasi memeriksa aturan kontrol dan kelengkapan;
  \item pelaporan mengumpulkan dan mengurutkan diagnostik.
\end{enumerate}

\begin{figure}[H]
\centering
\begin{tikzpicture}[node distance=6mm and 9mm]
  \node[phase, text width=2.4cm] (parse) {Parser};
  \node[phase, text width=2.4cm, right=of parse] (sem) {Type checker};
  \node[phase, text width=2.4cm, right=of sem] (val) {Validator};
  \node[phase, text width=2.4cm, below=of val] (report) {Diagnostik};
  \draw[arrow] (parse) -- (sem);
  \draw[arrow] (sem) -- (val);
  \draw[arrow] (val) -- (report);
\end{tikzpicture}
\caption{Pipeline validasi multi-pass}
\label{fig:ch16-pipeline}
\end{figure}

\subsection{Implementasi Validator}

Validator menggunakan traversal dengan konteks. Potongan fungsi
\code{visit} ditunjukkan pada \lstref{lst:ch16-visit}.

\includecode[
  language=C,
  firstline=31,
  lastline=90,
  firstnumber=31,
  caption={Fungsi \code{visit} pada \file{validate.c}},
  label={lst:ch16-visit}
]{../../code/subsetc/src/validate.c}

Setiap simpul fungsi memulai konteks baru dengan
\code{return\_type} dan \code{loop\_depth} yang sesuai. Simpul loop
menaikkan kedalaman sebelum mengunjungi anak dan menurunkannya
setelah selesai. Simpul \code{break} dan \code{continue} memeriksa
kedalaman loop.

\subsection{Validasi Fungsi}

Fungsi \code{validate\_program} menjadi titik masuk. Potongannya
ditunjukkan pada \lstref{lst:ch16-validate}.

\includecode[
  language=C,
  firstline=92,
  lastline=104,
  firstnumber=92,
  caption={Fungsi \code{validate\_program}},
  label={lst:ch16-validate}
]{../../code/subsetc/src/validate.c}

Validator mengembalikan jumlah kesalahan. Driver menggabungkan hasil
parser, type checker, dan validator untuk menentukan status akhir.
